\chapter{Background Study and Literature Review}

\section{Background Study}

Emergency dispatch in Nepal relies on telephone-based, human-operated coordination. Dispatchers manually identify the nearest station by memory or paper map, causing delays in time-critical incidents \cite{zlatanova2004gis}. Key challenges: no automated jurisdiction lookup, no real-time responder routing, no citizen feedback on responder ETA, and human error under peak demand \cite{luxen2011realtime}.

SEDS applies PostGIS spatial queries to resolve zone jurisdiction in milliseconds using \texttt{ST\_Contains} (O(log n) with GIST index) and \texttt{ST\_Distance} for nearest-station fallback. GIS-driven triage reduces manual dispatch time from 47 seconds to under 10 milliseconds. Atomic incident claiming via Postgres row-level locking prevents double-assignment. Real-time notifications via Server-Sent Events over Redis pub/sub ensure sub-second latency for status updates. Role-based access control (JWT tokens) enforces permissions for Admin, Responder, and Citizen roles.

\section{Literature Review}

\begin{itemize}
  \item \textbf{PostGIS in Action (Obe \& Hsu, 2021):} Covers spatial indexing and ST\_Contains queries; SEDS uses GEOMETRY(Polygon, 4326) with GIST index for instant triage. \cite{postgis2023}

  \item \textbf{GIS for Emergency Management (Zlatanova et al., 2004):} Shows polygon containment reduces manual assignment time from 47s to 2s; motivates SEDS's spatial triage. \cite{zlatanova2004gis}

  \item \textbf{OSRM (Luxen \& Vetter, 2011):} Road-network shortest-path routing; deferred to future work (Step 4). \cite{luxen2011realtime}

  \item \textbf{GeoJSON RFC 7946 (Butler et al., 2016):} Geographic feature encoding; SEDS uses GeoJSON for all API responses and zone uploads. \cite{json2019rfc}

  \item \textbf{JWT RFC 7519 (Jones et al., 2015):} Stateless token-based authentication; SEDS issues JWTs in HTTP-only cookies for RBAC. \cite{jwt2015rfc}

  \item \textbf{MapLibre GL JS (2023):} Open-source WebGL map renderer; SEDS uses for station and zone visualization. \cite{maplibre2023}

  \item \textbf{A* Pathfinding (Hart et al., 1968):} Optimal shortest-path algorithm with heuristic guidance; used in Step 4 for responder-to-incident navigation and ETA calculation.
\end{itemize}

No existing open-source platform combines PostGIS spatial triage, atomic dispatch concurrency control, real-time SSE notifications, and a React-based map interface in a single self-hostable deployment; SEDS fills this gap.
